% answers_short.tex —— 2004 年·全国卷Ⅳ·理科 · 简短答案速查
% 由 tools/make_short_answers.py 自 answers.tex 的 \gksolhead/\gkanswer 提取，
% 并与 2004/2004_数学真题_全国4理_答案要点.md 逐题比对后生成，勿手改（改源 MD 后重跑）。

\begin{gknotice}
\textbf{使用说明}：本节只给最终答案便于快速核对，完整推导见解析版。答案为编写者独立推导，\textbf{非官方标准答案}。
\end{gknotice}

\gkpart{选择题}
\begin{center}\begin{tabular}{ccccccccccccc}
\toprule
题号 & 1 & 2 & 3 & 4 & 5 & 6 & 7 & 8 & 9 & 10 & 11 & 12 \\
\midrule
答案 & D & C & A & D & A & B & C & A & B & A & B & C \\
\bottomrule
\end{tabular}\end{center}

\gkpart{填空题}
\begin{itemize}[leftmargin=2.6em]
  \item[\textbf{13.}] $28$
  \item[\textbf{14.}] $-\dfrac{1}{2}$
  \item[\textbf{15.}] $\dfrac{3}{4}$
  \item[\textbf{16.}] $2$
\end{itemize}

\gkpart{解答题}
\begin{itemize}[leftmargin=2.6em]
  \item[\textbf{17.}] $-\sqrt{2}$
  \item[\textbf{18.}] 在区间 $[0,2]$ 上的最大值为 $\ln 2-\dfrac{1}{4}$，最小值为 0
  \item[\textbf{19.}] （1）$\xi$ 的概率分布列为：
  \item[\textbf{20.}] （1）$96$；（2）证明见解析
  \item[\textbf{21.}] $\left[\dfrac{\sqrt{5}}{2}, \sqrt{5}\right]$
  \item[\textbf{22.}] （1）证明见解析；（2）$-\dfrac{\pi\mathrm{e}^\pi}{(\mathrm{e}^\pi+1)^2}$（或 $-\dfrac{\pi\mathrm{e}^{-\pi}}{(1+\mathrm{e}^{-\pi})^2}$）
\end{itemize}
